\subsection{含谱范数近似误差的完整有限步定理}\label{sec:inexact-theorem}
为使近似正交化进入主结论，以下将四项条件分别列明：目标有有限下界；满足式\eqref{eq:smooth}的 Frobenius 光滑性；随机梯度满足式\eqref{eq:noise}；实际方向与理想极因子的谱范数距离有统一上界。前三项沿用前文假设，第四项为
\begin{equation}\label{eq:delta-assumption}
 O_t^\star=\polar(M_t),\qquad
 \|O_t-O_t^\star\|_2\leq\delta\leq\bar\delta<1,
 \quad\text{对全部 }t\text{ 几乎处处成立}.
\end{equation}
这里 $O_t$ 是未作形状缩放的近似方向，更新为 $W_{t+1}=W_t-\eta O_t$；动量仍采用算法\ref{alg:exact}的初始化和普通 EMA。该条件定理不同时包含权重衰减、Nesterov 修正或辅助 AdamW。

\begin{theorem}[带谱范数近似误差的非渐近界]\label{thm:inexact}
在上述四项条件下，令 $\alpha=1-\beta$，则
\begin{equation}\label{eq:inexact-finite}
\boxed{\begin{aligned}
 A_T\leq\frac1{1-\delta}\biggl[&\frac\Delta{\eta T}
 +\frac{Lr\eta(1+\delta)^2}{2}\\
 &+2\sqrt r\left(\frac\sigma{T\alpha\sqrt B}
 +\sigma\sqrt{\frac\alpha B}
 +\frac{\beta L\eta\sqrt r(1+\delta)}\alpha\right)\biggr].
\end{aligned}}
\end{equation}
固定 $B$，取 $\eta=cT^{-3/4}$、$\alpha=T^{-1/2}$，且 $\delta\leq\bar\delta<1$ 与 $T$ 无关，则 $A_T=O(T^{-1/4})$。特别地，$\delta=O(T^{-1/4})$ 并最终小于 $\bar\delta$ 也是充分条件。
\end{theorem}
\begin{proof}
设 $R_t=O_t-O_t^\star$。由谱范数与核范数对偶关系，
\begin{align}
 \inner{H_t}{O_t}
 &=\inner{H_t}{O_t^\star}+\inner{H_t}{R_t}\notag\\
 &\geq(1-\delta)\nuc{H_t}-2\sqrt r\Frob{E_t}.
\end{align}
并且 $\|O_t\|_2\leq1+\delta$，所以 $\|O_t\|_F\leq\sqrt r(1+\delta)$。代入下降引理得
\begin{equation}\label{eq:inexact-descent}
 f(W_{t+1})\leq f(W_t)-\eta(1-\delta)\nuc{H_t}
 +2\eta\sqrt r\Frob{E_t}+\frac{Lr\eta^2(1+\delta)^2}{2}.
\end{equation}
式\eqref{eq:erecursion}的代数递推保持不变；噪声和初始化项也不变。实际更新的长度改变了漂移界：
\begin{equation}
 \|H_{t-1}-H_t\|_F\leq L\eta\sqrt r(1+\delta).
\end{equation}
因此式\eqref{eq:trackingavg}的最后一项变为 $\beta L\eta\sqrt r(1+\delta)/\alpha$。将这一平均误差界代入式\eqref{eq:inexact-descent}的望远镜求和，利用目标下界，再除以 $\eta T(1-\delta)$，即得式\eqref{eq:inexact-finite}。代入所给调度后，初始目标差项、噪声项和漂移项均不超过常数乘 $T^{-1/4}$，其余项阶数更高；$1/(1-\delta)\leq1/(1-\bar\delta)$，所以得到所述速率。
\end{proof}

\paragraph{为什么不必强制 $\delta\to0$？}
谱范数误差通过 $|\langle H_t,R_t\rangle|\leq\delta\|H_t\|_*$ 相对吸收进下降项。因此统一的 $\delta<1$ 已是这条证明路线的充分条件，而不是必须额外假定 $\delta=O(T^{-1/4})$。当 $\delta=0$ 时，式\eqref{eq:inexact-finite}退化为精确版本的界。若 $\delta\geq1$，该证明失去正下降系数；这只表示本充分界不可用，并不说明实际算法必然发散。

\subsection{谱范数光滑性、固定步数与适用范围}
若采用谱范数几何，梯度的自然对偶范数为核范数。一个明确的充分光滑条件是
\begin{equation}\label{eq:spectral-smooth}
 \|\nabla f(X)-\nabla f(Y)\|_*\leq L_s\|X-Y\|_2.
\end{equation}
它推出以 $L_s\|Y-X\|_2^2/2$ 为二次项的下降引理。重复上述证明，曲率项不再含 $r$，而梯度漂移用 $\|H_t-H_{t-1}\|_F\leq\|H_t-H_{t-1}\|_*\leq L_s\eta(1+\delta)$ 控制，得到明确的谱范数几何版本：
\begin{equation}\label{eq:spectral-inexact-finite}
\begin{aligned}
 A_T\leq\frac1{1-\delta}\biggl[&\frac\Delta{\eta T}+\frac{L_s\eta(1+\delta)^2}{2}\\
 &+2\sqrt r\left(\frac\sigma{T\alpha\sqrt B}+\sigma\sqrt{\frac\alpha B}
 +\frac{\beta L_s\eta(1+\delta)}\alpha\right)\biggr].
\end{aligned}
\end{equation}
外部的 $2\sqrt r$ 仍由方差所用的 Frobenius 范数转换而来。同样代入 $\eta=cT^{-3/4}$、$\alpha=T^{-1/2}$，得到平均核范数梯度的 $O(T^{-1/4})$ 界。不能只写“谱范数光滑”而不说明梯度在哪个对偶范数下计量。Frobenius 光滑可给出保守常数 $L_s=rL$，所以这两种写法不会自动带来维度上的改进。

Higham\citep{higham2008}讨论了极分解与矩阵函数迭代；本文五次多项式的具体系数取自 Muon 的工程实现\citep{jordan2024}，两者不应混作同一收敛保证。对固定 $K$，合成标量多项式 $\phi^{\circ K}(s)$ 在 $s\downarrow0$ 时趋于零，而理想极因子在非零支撑上的奇异值是 1。因此，在任意病态谱上不能仅由“执行五步”推出一个与输入无关的 $\bar\delta<1$。若要把 $\delta$ 的预算调度落实为 $K$ 的调度，还需要谱条件与多项式误差分析。

更直接地，经验五次多项式满足
\begin{equation}
 \phi(1)=3.4445-4.7750+2.0315=0.7010\ne1.
\end{equation}
所以 1 不是这个固定多项式的不动点。由连续性，反复使用同一组系数的奇异值序列不能以 1 为极限；单纯增大 $K$ 不会自动使谱范数误差趋于零。若需要可趋零的精确极因子误差保证，就必须采用具有相应保证的近似器或系数方案，再分析其谱条件。本文的条件定理不把这种实现保证作为已经证明的事实。

新增实验会在固定检查点记录 $\delta$ 的观测值，但稀疏检查点不能证明式\eqref{eq:delta-assumption}对全程成立。实际 ResNet-18 使用固定 $K$、Nesterov、形状缩放和权重衰减；其结果是实用配置的经验比较。精确版本、条件近似版本和实用混合版本的假设边界分别保留，不将一个版本的结论替换成另一个版本的保证。
